\section{Discussion}

The comparison against \texttt{opt -O1} is the honest heart of this report. On
\texttt{dead\_code}, \texttt{opt -O1} reduces the function far below what the
framework does, because it recognizes that the entire computation is dead and
removes the loop, while the framework's block local dead code elimination only
removes the straight line dead chain. On \texttt{branchy} and \texttt{clean},
\texttt{opt -O1} restructures control flow and unrolls loops, which can raise the
static instruction count even as it lowers runtime; the framework does neither,
so its instruction count only ever goes down. These are not shortcomings to
apologize for. They are the visible boundary of a deliberately small pass set,
and drawing that boundary clearly was a goal.

Where the framework does have real redundancy to work with, on
\texttt{arith\_redundant}, it produces both a structural reduction and a runtime
improvement that exceeds the measurement noise. That is the case the passes were
built for, and it shows that the pieces compose as intended: canonicalization
feeding common subexpression elimination, simplification feeding dead code
elimination, iterated to a fixed point.

The most valuable outcome was not any single rewrite but the discipline the
project enforced. The differential checksum gate caught the class of bug that
matters most, a miscompile, before it could ever be mistaken for a good result.
The invalidation test made the pass manager preservation contract concrete rather
than assumed. And generating every figure and number from the recorded CSV means
the report cannot drift from the measurement.
